\documentclass[11pt,a4paper]{article}
\usepackage[margin=21mm,headheight=14pt,includeheadfoot]{geometry}
\usepackage{amsmath,booktabs,enumitem,fancyhdr}

\setlength{\parindent}{0pt}
\setlength{\parskip}{3pt}
\setlist[enumerate]{label=\textbf{(\alph*)},leftmargin=*,itemsep=3pt,topsep=3pt,parsep=0pt}
\pagestyle{fancy}
\fancyhf{}
\fancyhead[L]{\small COMP2040 --- Operating Systems}
\fancyhead[R]{\small Checkpoint 1}
\fancyfoot[C]{\small\thepage}
\newcommand{\code}[1]{\texttt{#1}}

\begin{document}
\begin{center}
  {\Large\bfseries COMP2040: Operating Systems\par}
  \vspace{3pt}
  {\large Checkpoint 1: Progress Report\par}
  \vspace{5pt}
  Nguyen Vu Nhat Minh \quad Student ID: V202401314\\
  {\small College of Engineering and Computer Science, VinUniversity}
\end{center}

This report presents my completed contribution as Person A: SJF, SRTF
and preemptive priority scheduling in \code{ossim}, together with their
hand checks and regression tests. FCFS and RR were supplied in the scaffold.

\section{Scheduler Data Structures and Design}

Each scheduler owns a separate \code{queue\_t}, an array-backed FIFO of
PCB pointers, and uses the existing \code{scheduler\_t} hooks. READY
processes join the tail. \code{q\_pop\_min} selects the smallest key;
equal keys retain ready-queue order, rather than PID order.

\begin{enumerate}
\item \textbf{SJF: non-preemptive shortest job first.}
The key, \code{proc\_burst\_remaining}, measures the \emph{current}
CPU burst, not all future CPU work. SJF never preempts. The process runs
until blocking or termination; consecutive \code{cpu} operations retain
the CPU through the kernel's \code{continuing} path.

\item \textbf{SRTF: shortest remaining time first.}
SRTF uses the same burst key. Before a CPU tick, it preempts if
\code{q\_min\_key} is \emph{strictly smaller} than the runner's remaining
burst; equality keeps the runner. Events are delivered first, so arrivals
and I/O completions can preempt in their own tick. Preempted processes
rejoin the tail.

\item \textbf{PRIO: preemptive priority with round-robin ties.}
The key is \code{priority}; a smaller number runs first. A strictly
higher priority preempts immediately. Equal priorities rotate after
\code{cfg.quantum} CPU ticks, using FIFO ties and tail reinsertion.
\code{on\_tick} increments \code{quantum\_used}; \code{on\_dispatch}
resets it, giving a fresh quantum after interruption. With only lower
priorities ready, expiry resets the quantum and the runner continues.
Consecutive CPU operations retain their quantum without redispatch.
\end{enumerate}

\textbf{Cost and prediction.}
Selection, removal and preemption scans cost \(O(n)\). With at most 64
processes, simple queues are adequate. \code{ossim} knows burst lengths
exactly; a real OS could predict the next burst with exponential averaging:
\[
  \tau_{n+1}=\alpha t_n+(1-\alpha)\tau_n,
  \qquad 0\leq\alpha\leq1,
\]
Here \(t_n\) is the last measured burst and \(\tau_n\) its estimate.
Choosing an initial estimate and \(\alpha\) would be necessary; prediction
is not implemented.

\newpage
\section{How the Schedulers Were Checked}

I traced \code{workloads/arrivals.txt} by hand before comparing with
\code{ossim}, using \code{--no-paging} and quantum 2 for PRIO.
Waiting counts READY ticks; turnaround is finish minus arrival, and
response is first dispatch minus arrival. Switches include the initial
dispatch.

\begin{center}
{\small
\begin{tabular}{@{}lcccc@{}}
\toprule
 & \multicolumn{3}{c}{Average time in ticks: hand / simulator} & Switches \\
\cmidrule(lr){2-4}
Scheduler & Waiting & Response & Turnaround & hand / simulator \\
\midrule
SJF & 10.00 / 10.00 & 9.00 / 9.00 & 17.83 / 17.83 & 8 / 8 \\
SRTF & 9.83 / 9.83 & 9.00 / 9.00 & 17.67 / 17.67 & 9 / 9 \\
PRIO & 12.83 / 12.83 & 8.33 / 8.33 & 20.67 / 20.67 & 12 / 12 \\
\bottomrule
\end{tabular}}
\end{center}

Each metric is checked in \code{tests/run.sh}; arithmetic and Gantt charts
are in \code{docs/checks/sjf.md}, \code{srtf.md} and \code{prio.md}.

\textbf{Regression and integration checks.}
SJF gives average waiting 3.00 for textbook bursts 24, 3, 3.
\code{sched\_edges.txt} covers FIFO ties, same-tick events and CPU
continuation; \code{prio\_check.txt} checks fresh quanta after interruption.
Producer/consumer and mixed paging workloads finish under all three
new schedulers, including mixed workloads with one frame.

GCC 15.2 on Ubuntu WSL built the code with
\code{-std=c99 -Wall -Wextra -Werror}. The recorded terminal summary is:
{\small
\begin{verbatim}
$ make test
...
54 passed, 5 skipped (not implemented yet), 0 failed
\end{verbatim}}
Another 30 runs completed ten non-deadlocking workloads under the three
new schedulers. The terminal evidence is saved in
\code{docs/checks/person\_a\_terminal.txt}.

\section{Contributions and Current Status}

\textbf{Nguyen Vu Nhat Minh (Person A, \code{garan1610}):}
implemented SJF/SRTF/PRIO; prepared hand checks, workloads, tests and
documentation; reviewed deadlock cases; fixed WSL shell line endings;
and strengthened \code{check}/\code{check\_no} to reject failed commands.

\textbf{Person B (\code{TheNameisSea}):}
implemented resource deadlock detection and its associated workloads,
tests and hand checks.

MLFQ and its review remain pending, accounting for one SKIP; the other
four are Checkpoint 2 checks. Teammate approval is pending. PRIO can
starve lower priorities without aging. Only GCC was verified because
Clang was unavailable. Codex assisted implementation, hand checks and
debugging.

\end{document}
